% Section 6: System
% Target length: 500-800 words
% Status: OUTLINE ONLY - grounded in actual current skill behavior (verified by
% reading SKILL.md source directly -- see notes below). Needs prose + the worked
% table turned into a proper LaTeX table/figure.

\section{System}
\label{sec:system}

\subsection{Current Pipeline}
% TODO: expand into prose. Facts verified directly from source, cite paths as
% footnotes or in a reproducibility appendix rather than in-text:
\begin{itemize}
    \item \texttt{/book-chapter-generator}: assigns concepts from
          \texttt{learning-graph.json} to chapters. Verified behavior: uses
          topological sorting for validity and heuristic clustering of ``natural
          learning units''; does not currently use in-degree or CIS.
    \item \texttt{/chapter-content-generator}: generates chapter prose from a
          chapter's ``Concepts Covered'' list. Verified current instruction: flat
          target of \textbf{3000--5000 words per chapter} and \textbf{4-6 non-text
          elements per chapter}, independent of concept importance.
    \item Neither skill currently reads or computes in-degree, CIS, or any other
          graph-centrality measure. The subjectively-observed quality difference
          between Claude Sonnet 5 and Google Antigravity/Gemini-generated content
          (\S\ref{sec:introduction}) therefore cannot be attributed to an explicit
          mechanism in the shared tooling -- if real, it is either an emergent
          model behavior or an artifact of a different, unrecorded prompting
          session. See limitations (\S\ref{sec:discussion}) -- the original
          Gemini/Antigravity output was not preserved and could not be reanalyzed.
\end{itemize}

\subsection{Proposed Extension: CIS-Annotated Chapter Scaffold}
% TODO: expand. This is the concrete system contribution -- write as a diff
% against current behavior.
\begin{itemize}
    \item Compute CIS once per book, over the full learning graph (not
          per-chapter -- see global-vs-local normalization decision,
          \S\ref{sec:formal-definitions}).
    \item Extend \texttt{/book-chapter-generator}'s chapter scaffold output: each
          entry in ``Concepts Covered'' gets annotated with its CIS score, its
          Elaboration Score $E(c)$ (Definition 4, \S\ref{sec:formal-definitions}),
          and an assigned elaboration tier.
    \item Extend \texttt{/chapter-content-generator}'s instructions: replace the
          flat word-count instruction with a per-concept budget table (words +
          expected non-text elements), computed from the annotated scaffold.
\end{itemize}

\subsection{Worked Example: Algebra I, Chapter 1}

% TODO: expand surrounding prose. Real numbers, computed directly from
% docs/learning-graph/learning-graph.json in the Algebra I repo (200 concepts,
% 277 edges, confirmed valid DAG). Full ranked data:
% data/algebra1-concept-impact.csv. Chapter-1-specific budget table source:
% data/algebra1-ch1-elaboration-budget.csv.
%
% Chapter 1 ("Foundations of Algebra", 17 concepts) is the sharpest test case:
% it contains several concepts with LOW in-degree but very HIGH CIS due to
% transitive reach (Constant: indeg=1, rank #53 by indeg, but CIS=238, rank #4;
% Coefficient: indeg=2, rank #35, CIS=239, rank #3; Term: indeg=4, rank #20,
% CIS=237, rank #5). This directly demonstrates the failure mode of raw
% in-degree that motivates CIS.
%
% Normalization: GLOBAL, tiered by Elaboration Score E(c), NOT population
% percentile rank (resolved after the Chapter 12 validation check below found
% percentile rank was structurally broken -- see \S\ref{sec:formal-definitions}
% Definition 4). Tier cut points: E>=0.5 = Tier A, 0.2<=E<0.5 = Tier B, E<0.2 =
% Tier C. Word budget: fixed chapter total of 6800 words (matches the current,
% already-published Chapter 1 word count -- see
% \texttt{docs/learning-graph/chapter-metrics.md} in the Algebra I repo),
% redistributed by log-weighted proportional allocation using each concept's
% GLOBAL CIS value, floor of 150 words/concept -- this word-budget formula is
% unaffected by the percentile-rank bug fix, since it already used log(CIS+1)
% directly rather than population rank.

\begin{table}[htbp]
\centering
\caption{CIS-weighted elaboration budget for Algebra I, Chapter 1 (``Foundations of Algebra''), 17 concepts, global normalization by Elaboration Score $E(c)$, fixed 6{,}800-word chapter total.}
\label{tab:ch1-budget}
\begin{tabular}{lrrrl r}
\toprule
Concept & in-deg & CIS & $E(c)$ & Tier & Words \\
\midrule
Variable                & 8  & 619 & 1.000 & A & 630 \\
Number                   & 13 & 369 & 0.920 & A & 592 \\
Coefficient              & 2  & 239 & 0.852 & A & 559 \\
Constant                 & 1  & 238 & 0.852 & A & 559 \\
Term                     & 4  & 237 & 0.851 & A & 559 \\
Expression               & 4  & 133 & 0.762 & A & 516 \\
Equation                 & 5  & 111 & 0.734 & A & 503 \\
Monomial                 & 1  & 76  & 0.676 & A & 475 \\
Inequality               & 4  & 13  & 0.410 & B & 347 \\
Order of Operations      & 1  & 7   & 0.323 & B & 305 \\
Like Terms               & 1  & 7   & 0.323 & B & 305 \\
Evaluating Expressions   & 3  & 6   & 0.303 & B & 295 \\
Combining Like Terms     & 4  & 6   & 0.303 & B & 295 \\
Substitution              & 2  & 3   & 0.216 & B & 254 \\
Simplifying Expressions  & 0  & 1   & 0.108 & C & 202 \\
Expanding Expressions    & 0  & 1   & 0.108 & C & 202 \\
Prime Factorization      & 0  & 1   & 0.108 & C & 202 \\
\bottomrule
\end{tabular}
\end{table}

Chapter~1 yields 8 Tier-A, 6 Tier-B, and 3 Tier-C concepts -- still skewed
toward heavy treatment, as expected for the book's foundational chapter, but
no longer degenerate the way the earlier percentile-rank scheme was.

\subsection{Validation: Algebra I, Chapter 12}

% TODO: expand surrounding prose.
Chapter 1 alone cannot validate the tiering rule, since a heavy skew toward
Tier A is plausible there regardless of whether the method is working
correctly. As a second, contrasting check we computed the same table for
Chapter 12 (``Exponential Functions'', 10 concepts) -- a late, specialized
chapter where 8 of 10 concepts have in-degree 0 (true terminal concepts with
no downstream dependents). Word budget fixed at the chapter's current
published total, 3{,}722 words
(\texttt{docs/learning-graph/chapter-metrics.md}).

This check is what first exposed the population-percentile bug (Definition 4,
\S\ref{sec:formal-definitions}): under percentile rank, every one of these
zero-in-degree, CIS-floor concepts landed at the 51.5th percentile (since
103/200 concepts in the book, 51.5\%, sit at the same CIS floor of 1), safely
above the 35th-percentile Tier-C cutoff -- producing zero Tier-C concepts in
a chapter that should plausibly have several. The corrected Elaboration Score
$E(c)$ resolves this by scoring on the log-scaled CIS value rather than
population rank, so tied floor concepts separate cleanly from true hubs.

\begin{table}[htbp]
\centering
\caption{CIS-weighted elaboration budget for Algebra I, Chapter 12 (``Exponential Functions''), 10 concepts, global normalization by Elaboration Score $E(c)$, fixed 3{,}722-word chapter total.}
\label{tab:ch12-budget}
\begin{tabular}{lrrrl r}
\toprule
Concept & in-deg & CIS & $E(c)$ & Tier & Words \\
\midrule
Exponential Function              & 6 & 10 & 0.373 & B & 697 \\
Exponential Growth                & 2 & 3  & 0.216 & B & 466 \\
Exponential Decay                 & 1 & 2  & 0.171 & C & 401 \\
Growth Factor                     & 0 & 1  & 0.108 & C & 308 \\
Decay Factor                      & 0 & 1  & 0.108 & C & 308 \\
Initial Value                     & 0 & 1  & 0.108 & C & 308 \\
Exponential Models                & 0 & 1  & 0.108 & C & 308 \\
Comparing Linear and Exponential  & 0 & 1  & 0.108 & C & 308 \\
Graphing Exponentials             & 0 & 1  & 0.108 & C & 308 \\
Applications of Systems           & 0 & 1  & 0.108 & C & 308 \\
\bottomrule
\end{tabular}
\end{table}

Chapter~12 yields 0 Tier-A, 2 Tier-B, and 8 Tier-C concepts -- the mirror
image of Chapter~1's profile, as expected for a late, specialized chapter
built mostly from terminal concepts. The two chapters together are the
validation check this tiering rule needed: a real, opposite-direction split
between a foundational and a specialized chapter, rather than the same
degenerate all-A or all-B/no-C result regardless of chapter content. The
0.5/0.2 cut points for $E(c)$ are a first pass validated on exactly these two
chapters; treat them as a starting point to revisit with more chapters, not a
settled constant (\S\ref{sec:evaluation}).
